\documentclass[11pt]{article}
\usepackage[T1]{fontenc}
\usepackage{lmodern}
\usepackage[a4paper,margin=27mm]{geometry}
\usepackage{amsmath,amssymb,amsthm,mathtools}
\usepackage{microtype}
\usepackage[colorlinks=true,linkcolor=blue,citecolor=blue,urlcolor=blue]{hyperref}
\usepackage{enumitem}
\setlist{nosep}
\newtheorem{theorem}{Theorem}[section]
\newtheorem{proposition}[theorem]{Proposition}
\newtheorem{lemma}[theorem]{Lemma}
\newtheorem{corollary}[theorem]{Corollary}
\newtheorem{remark}[theorem]{Remark}
\newcommand{\PP}{\mathbb P}
\newcommand{\one}{\mathbf 1}
\newcommand{\Ctwo}{C_2}
\newcommand{\twin}{\mathcal T}
\newcommand{\cset}{\mathcal C}
\newcommand{\aset}{\mathcal A}
\newcommand{\phimin}{P^-}
\newcommand{\thetap}{\vartheta}
% Optional author details are separate from the project contact below.
\newcommand*{\ManuscriptAuthor}{}
\newcommand*{\ManuscriptEmail}{}
\title{A Long-Divisor Sign Classification on Three-Prime Products\\
\large and an Upper Sieve for Their Shift by Two}
\author{\ManuscriptAuthor\ifx\ManuscriptEmail\empty\else
  \ifx\ManuscriptAuthor\empty\else\\\fi
  {\small\href{mailto:\ManuscriptEmail}{\expandafter\nolinkurl\expandafter{\ManuscriptEmail}}}\fi}
\date{\small DEEPAK KUMAR, project coordinator and contact\\[2pt]
  \href{mailto:admin@cosmosgeek.org}{\nolinkurl{admin@cosmosgeek.org}}\\[4pt]
  Research draft --- 25 November 2026}
\begin{document}
\maketitle
\begin{abstract}
For a squarefree composite integer $M$ near $X$ we study the truncated
long-divisor convolution
\[
 F_D(M)=\sum_{\substack{k\mid M\\1<k<M/D}}\mu(M/k)\log k,
 \qquad D=X^{0.24995}.
\]
We determine $F_D$ exactly on products of two and three distinct primes.
On semiprimes it is negative.  On three-prime products its sole positive
region consists of $M=rsp$, where the two smaller primes satisfy
$r,s\le D<rs$, and the weight there is exactly $\log p$.
With the additional cutoff \(r>X^{1/20}\), the unshifted weighted mass of this region is
$(\mathcal I+o(1))X$, where $\mathcal I$ is an explicit dilogarithmic
constant, $0.611418619809\ldots$.
We then retain the genuine shift-two condition and prove a uniform
linear-sieve upper bound for the part on which $M+2$ is squarefree with
exactly two or three sufficiently large prime factors; on a fixed factor
box the normalized upper bound is $0.005998386331\ldots$.
A companion positive lower bound is established only for the larger
family in which $M+2$ is rough.  None of these results bounds the
complete signed remainder sufficiently to prove the twin-prime conjecture.
\end{abstract}

\section{Definitions and statements}

Throughout, $X\to\infty$, and a letter denoting a prime is always
explicitly so designated.  Write $\omega(n)$ for the number of
\emph{distinct} prime factors, $\mu$ for the M\"obius function,
$\phimin(n)$ for the least prime factor, and
$\thetap(n)=\one_{n\text{ prime}}\log n$.
We use
\[
 \delta=\frac1{20},\qquad t=\frac{4999}{20000}=0.24995,
 \qquad D=X^t,\qquad Y=X^{1/21}.
\]
(The two cutoff exponents are different: $\delta>1/21$.)
For squarefree $M>1$ put
\begin{equation}\label{eq:F}
 F_D(M):=\sum_{\substack{k\mid M\\1<k<M/D}}
                 \mu(M/k)\log k.
\end{equation}
The strict inequality in this definition fixes the endpoint convention.
We emphasize that both $k$ and $M/k$ are numerical divisors, not
necessarily primes.

\begin{theorem}[Exact low-rank classification]\label{th:classification}
Let $M\in(X,2X]$ be squarefree.  If $M=rs$ with primes $r<s$, then
\begin{equation}\label{eq:rank2}
 F_D(rs)=\begin{cases}-\log r,&r\le D,\\
                         -\log(rs),&r>D.
          \end{cases}
\end{equation}
If $M=rsp$ with primes $r<s<p$, then
\begin{equation}\label{eq:rank3}
 F_D(rsp)=\begin{cases}
  -\log(rsp),&r>D,\\
  -\log r,&r\le D<s,\\
  0,&r<s\le D\ \text{and }rs\le D,\\
  \log p,&r<s\le D\ \text{and }rs>D.
 \end{cases}
\end{equation}
In particular, the only positive rank-three sector compatible with
$\phimin(M)>X^\delta$ is
\begin{equation}\label{eq:C}
 \cset_X=\left\{M=rsp:\begin{array}{l}
 X<M\le2X,\ r,s,p\text{ distinct primes},\ r<s<p,\\
 X^\delta<r<s\le X^t,\ rs>X^t
 \end{array}\right\}.
\end{equation}
Every $M\in\cset_X$ has $p>X^{1-2t}=X^{0.50010}$, and on this set
$F_D(M)=\log p$.
\end{theorem}

Define
\begin{equation}\label{eq:I}
 \mathcal I=\frac12\int_\delta^t\frac{du}{u}
              \int_{\max(\delta,t-u)}^t\frac{dv}{v}.
\end{equation}
Let $\operatorname{Li}_2(z)=-\int_0^z\log(1-v)\,dv/v$ be the dilogarithm.

\begin{theorem}[The unshifted weighted three-prime mass]\label{th:unshifted}
The sum over the \emph{genuine} primes in \eqref{eq:C} satisfies
\begin{equation}\label{eq:unshifted}
 \sum_{M\in\cset_X}\log P^+(M)=(\mathcal I+o(1))X,
\end{equation}
where $P^+(M)=p$ is its largest prime factor.  With $a=\delta/t$,
\begin{equation}\label{eq:dilog}
 \mathcal I=\frac12\left\{
 \operatorname{Li}_2(1-a)-\operatorname{Li}_2(a)
 +\log(1/a)\log\frac1{1-a}\right\}
 =0.61141861980926243232\ldots .
\end{equation}
Requiring in addition $M\equiv2\pmod3$ changes the main term in
\eqref{eq:unshifted} to $(\mathcal I/2+o(1))X$.
\end{theorem}

The unshifted mass in Theorem \ref{th:unshifted} is of order $X$.
The shifted estimates below use a normalization of order $X/\log X$.
We prove a lower bound of this order for the rough family defined below.
The more restrictive weighted sum has only an upper bound here.
To define the latter, put
\[
 C_2=\prod_{\substack{q>2\\q\text{ prime}}}
          \left(1-\frac1{(q-1)^2}\right),\quad
 J_X=\int_X^{2X}\frac{du}{\log u},\quad B_X=2C_2J_X
    \sim\frac{2C_2X}{\log X}.
\]
For squarefree $N>1$ define the quadratic truncated divisor weight
\begin{equation}\label{eq:G}
 G_2(N)=\sum_{d\mid N}\mu(d)
                  \left(1-2\frac{\log d}{\log N}\right)_+^2,
 \qquad v_+=\max(0,v),
\end{equation}
and the prescribed neighboring-integer gate
\begin{equation}\label{eq:A}
 \aset_X=\{N:\ N\text{ squarefree},\quad
 \omega(N)\in\{2,3\},\quad \phimin(N)>Y\}.
\end{equation}
The genuine coupled contribution from the positive sector is
\begin{equation}\label{eq:S}
 S_{+,3}(X)=\sum_{M=rsp\in\cset_X}
       (\log p)G_2(rsp+2)\one_{\{rsp+2\in\aset_X\}}.
\end{equation}
For $0<a_1<a_2<b_1<b_2<t$, with $a_1+b_1>t$ and
$a_1>\delta$, denote by $S_{+,3}^{\Box}(X)$ the same sum with
\[
 X^{a_1}<r\le X^{a_2},\qquad
 X^{b_1}<s\le X^{b_2}.
\]
This rectangle lies entirely in \eqref{eq:C}.

\begin{lemma}[Positivity and size of the gate weight]\label{lem:G}
For every $N\in\aset_X$, one has $0\le G_2(N)\le1$.
Writing the prime logarithm shares as nonnegative numbers adding to one,
$G_2(rs)=4x(1-x)$ for two distinct factors.  For three factors with
shares $x,y,z$, one has
\begin{equation}\label{eq:G3}
 G_2(rsp)=\begin{cases}
  2-4(x^2+y^2+z^2),&\max(x,y,z)\le\frac12,\\
  8xy,& z>\frac12,\ \text{where }x+y=1-z.
 \end{cases}
\end{equation}
\end{lemma}

\begin{theorem}[A genuine shifted upper sieve]\label{th:shifted}
Let $a_1,a_2,b_1,b_2$ obey the preceding strict inequalities.  Then
\begin{equation}\label{eq:boxbound}
 0\le\limsup_{X\to\infty}\frac{S_{+,3}^{\Box}(X)}{B_X}
 \le14\log\frac{a_2}{a_1}\log\frac{b_2}{b_1}.
\end{equation}
For the full sector \eqref{eq:C},
\begin{equation}\label{eq:wholebound}
 0\le\limsup_{X\to\infty}\frac{S_{+,3}(X)}{B_X}
 \le14\mathcal I=8.5598606773296740\ldots .
\end{equation}
For the explicit box
\begin{equation}\label{eq:box}
 (a_1,a_2,b_1,b_2)=(0.125,0.1305,0.200,0.202)
\end{equation}
this becomes the concrete normalized bound
\begin{equation}\label{eq:numericbox}
 \limsup_{X\to\infty}\frac{S_{+,3}^{\Box}(X)}{B_X}
 \le0.005998386330999321\ldots .
\end{equation}
These are bounds for the original \emph{coupled} weight in \eqref{eq:S}.
They do not assert an asymptotic for $S_{+,3}$.
\end{theorem}

\begin{theorem}[The larger rough shifted family]\label{th:rough}
Let $\mathcal R_X^{\Box}$ be the sum obtained by replacing the
factor gate and $G_2$ in $S_{+,3}^{\Box}$ with
$\one_{\{\phimin(rsp+2)>Y\}}$ alone.  Then
\begin{align}\label{eq:rough}
 14(\log2)\log\frac{a_2}{a_1}\log\frac{b_2}{b_1}
 &\le \liminf_{X\to\infty}\frac{\mathcal R_X^{\Box}}{B_X}
 \le\limsup_{X\to\infty}\frac{\mathcal R_X^{\Box}}{B_X}\\
 &\le14\log\frac{a_2}{a_1}\log\frac{b_2}{b_1}.
 \notag
\end{align}
In particular, the larger rough family has positive weighted mass
of order $X/\log X$ on every such fixed nonempty box.  It can contain
neighbors with more than three prime factors and therefore does not
supply a lower bound for \eqref{eq:S}.
\end{theorem}

\section{The exact signed convolution}

\begin{proof}[Proof of Theorem \ref{th:classification}]
For any positive integer $n$ the standard Dirichlet-convolution
identity $\mu*\log=\Lambda$ reads
\begin{equation}\label{eq:mob}
 \sum_{d\mid n}\mu(d)\log(n/d)=\Lambda(n).
\end{equation}
If $M$ is squarefree with at least two distinct prime factors,
$\Lambda(M)=0$.  In the definition of $F_D$ put $d=M/k$.
The window $1<k<M/D$ becomes $D<d<M$.  The $d=M$ term in
\eqref{eq:mob} is zero because it contains $\log1$.  Therefore
\begin{equation}\label{eq:cut}
 F_D(M)=-\sum_{\substack{d\mid M\\d\le D}}
                 \mu(d)\log(M/d).
\end{equation}
This formula includes $d=1$ and all eligible nontrivial divisors.
The $d=1$ term is $-\log M$; omitting it would change the signs.

If $M=rs$ with $r<s$, then $s>D$, since $rs>X$ but
$D^2=X^{0.49990}<X$.  The only divisors up to $D$ are $1$ and,
when $r\le D$, $r$.  Thus \eqref{eq:cut} gives \eqref{eq:rank2}.

For $M=rsp$, at least one prime factor exceeds $D$, since
$D^3=X^{0.74985}<X<M$.  The largest prime $p$ exceeds $D$.
If $r>D$, only $d=1$ is available.  If $r\le D<s$, the
contribution $+\log(sp)$ from $d=r$ cancels all but $-\log r$.
If $r<s\le D$, the contributions $d=1,r,s$ total
\[
 -\log(rsp)+\log(sp)+\log(rp)=\log p.
\]
When $rs\le D$, the further contribution of $d=rs$ is
$-\log p$ and the total is zero.  When $rs>D$ it remains $\log p$.
This proves the exhaustive four-case formula.
On the positive case $rs\le D^2$ and $M>X$, whence
$p>X/D^2=X^{1-2t}$, strictly greater than $s\le D$.
Imposing the stipulated floor $\phimin(M)>X^\delta$ gives
\eqref{eq:C} without altering any of the signs.
\end{proof}

\begin{proof}[Proof of Lemma \ref{lem:G}]
For a squarefree semiprime the two nontrivial singleton terms of
\eqref{eq:G} have arguments $x$ and $1-x$, and the $d=N$ term
vanishes.  Only the term for the smaller share can be nonzero,
so $G_2=1-(1-2x)^2=4x(1-x)$.
For three prime factors, the subset of size two corresponding to
$x+y$ contributes $(2z-1)_+^2$ because $x+y=1-z$.
If all shares are at most $1/2$, all size-two terms vanish and
\[
 G_2=1-\sum_{w\in\{x,y,z\}}(1-2w)^2
      =2-4(x^2+y^2+z^2).
\]
Cauchy--Schwarz gives $x^2+y^2+z^2\ge1/3$,
while $x^2+y^2+z^2\le(\max w)\sum w\le1/2$.
Thus this value lies in $[0,2/3]$.
If $z>1/2$, then only the singleton terms for $x,y$ and the
doubleton term for $x+y$ remain, and expansion gives $8xy$.
It is between zero and $2(x+y)^2<1/2$.
The stated uniform bound follows.
\end{proof}

\section{The ordinary mass and its logarithmic geometry}

\begin{proof}[Proof of Theorem \ref{th:unshifted}]
Fix two primes $r<s$ satisfying the two small-prime constraints
in \eqref{eq:C}.  Since $rs\le D^2=X^{0.49990}$,
$Z=X/(rs)\ge X^{0.50010}$, and the prime number theorem gives,
uniformly for all such pairs,
\begin{equation}\label{eq:pnt}
 \sum_{Z<p\le2Z}\log p=Z+O_A(Z(\log X)^{-A})
\end{equation}
for every fixed $A>0$.  Here $p>s$ automatically.  Summation
of \eqref{eq:pnt} over $r,s$ therefore yields
\begin{equation}\label{eq:primeharmonic}
 \sum_{M\in\cset_X}\log p
  =X\sum_{\substack{X^\delta<r<s\le X^t\\rs>X^t}}
       \frac1{rs}+o(X),
\end{equation}
using $\sum_{X^\delta<q\le X^t,q\text{ prime}}q^{-1}=O(1)$.
By Mertens' reciprocal-prime estimate, the measure on the logarithmic
coordinate $u=\log q/\log X$ converges on each fixed positive-power
interval to $du/u$.  The boundary $u+v=t$ and the diagonal $u=v$
have zero limiting measure.  Thus the harmonic sum on the right
of \eqref{eq:primeharmonic} tends to
\[
 \frac12\int_\delta^t\frac{du}{u}
       \int_{\max(\delta,t-u)}^t\frac{dv}{v}=\mathcal I.
\]
The factor $1/2$ accounts for the unique order $r<s$ in the symmetric
full square, not for any selection of two among three prime factors.

Set $a=\delta/t$, so $0<a<1/2$, and scale both integration variables
by $t$.  Splitting the first variable at $1-a$ gives
\[
 2\mathcal I=
 \int_a^{1-a}\frac{-\log(1-u)}u\,du
 +\log(1/a)\int_{1-a}^1\frac{du}{u}.
\]
Integrating the first term using the definition of $\operatorname{Li}_2$
and the second term directly proves \eqref{eq:dilog}.
For the congruence $M\equiv2\pmod3$, both $r,s$ are eventually
coprime to $3$.  It prescribes one of the two reduced prime residue
classes $p\pmod3$.  The prime number theorem in the fixed modulus $3$
replaces \eqref{eq:pnt} by $Z/2+o(Z)$ uniformly, giving half the main.
No condition on the prime factors of $M+2$ has entered this proof.
\end{proof}

\begin{remark}[An optional unweighted companion]
The same argument with the unweighted prime number theorem shows
\[
 \#\cset_X\sim\frac{X}{2\log X}
 \int_\delta^t\frac{du}{u}
 \int_{\max(\delta,t-u)}^t\frac{dv}{v(1-u-v)}.
\]
The factor $(1-u-v)^{-1}$ accounts for the prime density at
$p\asymp X^{1-u-v}$.  It is absent from the \emph{logarithmically
weighted} main constant $\mathcal I$.
\end{remark}

\section{A shifted upper and lower linear sieve}

The analytic input here is the standard one-dimensional linear
sieve, for which the needed statements and the functions $f,F$ are
recorded in \cite[Lemma~2.5]{MZA} (see also
\cite{OperaDeCribro}).  We give the reduction in full to identify
the exact shifted residues and the uniformity required by that result.

\begin{lemma}[Prime-variable roughness sieve]\label{lem:sieve}
Uniformly for $m=rs$ with distinct primes
$X^\delta<r<s\le D$, and $Z=X/m$, let
\[
 \mathcal R_m(X)=\sum_{Z<p\le2Z}\thetap(p)
                    \one_{\{\phimin(mp+2)>Y\}}.
\]
Then, as $X\to\infty$,
\begin{equation}\label{eq:permbound}
 \left(\frac{4C_2\log2}{3}+o(1)\right)\frac{X}{m\log Y}
 \le \mathcal R_m(X)
 \le\left(\frac{4C_2}{3}+o(1)\right)\frac{X}{m\log Y}.
\end{equation}
The $o(1)$ terms are uniform over all the prescribed $m$.
\end{lemma}

\begin{proof}
The primes $r,s$ both exceed $Y$, since $\delta>1/21$.
For every odd squarefree integer $d$ composed of primes at most $Y$,
$(m,d)=1$ and
\begin{equation}\label{eq:congruence}
 d\mid(mp+2)\quad\Longleftrightarrow\quad
 p\equiv-2\overline m\pmod d,
\end{equation}
a reduced residue.  The prime $2$ never divides $mp+2$ on
this family since $m,p$ are odd.  For such $d$ set
\[
 A_m(d)=\sum_{\substack{Z<p\le2Z\\p\equiv-2\overline m\;(\bmod d)}}
        \log p.
\]
The expected main is $Z/\varphi(d)$, so the multiplicative local
sieve density is $g(q)/q=1/(q-1)$ for odd $q\le Y$, with $g(2)=0$.
Moreover
\[
 Z\ge X^{1-2t}=X^{0.50010},\qquad
 Y^3=X^{1/7}\le Z^{1/2-\epsilon_0}
\]
for some fixed $\epsilon_0>0$.  The maximal
Bombieri--Vinogradov theorem, with its maximum over reduced residues
and interval endpoints, gives for every fixed $A>0$
\begin{equation}\label{eq:BV}
 \sum_{\substack{d\le Y^3\\d\text{ odd}}}
 \max_{(a,d)=1}
 \left|\sum_{\substack{Z<p\le2Z\\p\equiv a\;(\bmod d)}}
                 \log p-\frac Z{\varphi(d)}\right|
 \ll_A\frac{Z}{(\log X)^A}.
\end{equation}
The inequality is uniform in $m$, because the maximum already includes
the changing residue $-2\overline m$.  Any \emph{fixed} divisor weight
needed by the sieve can also be included in \eqref{eq:BV} at the cost
of choosing a larger $A$: apply Cauchy--Schwarz to the original
absolute discrepancy sum, use the elementary bound
$\sum_{d\le Y^3}\tau(d)^K\ll_K Y^3(\log X)^{O_K(1)}$,
and use $|A_m(d)-Z/\varphi(d)|\ll (Z/d+1)\log(2Z)$.

The sifting density is
\begin{align}\label{eq:prod}
 V(Y)&=\prod_{\substack{3\le q\le Y\\q\text{ prime}}}
           \left(1-\frac1{q-1}\right)\\
 &=\prod_{3\le q\le Y}\left(1-\frac1q\right)
                  \left(1-\frac1{(q-1)^2}\right)
  =\frac{2C_2e^{-\gamma}+o(1)}{\log Y}.
 \notag
\end{align}
This uses Mertens' product formula and the absolute convergence
of the twin-prime Euler factor.  The upper and lower linear sieve
with level $Y^3$ has parameter $s=\log(Y^3)/\log Y=3$.
Its standard functions satisfy
\[
 F(3)=\frac{2e^\gamma}{3},\qquad
 f(3)=\frac{2e^\gamma}{3}\log2.
\]
The sieve thus bounds $\mathcal R_m$ between
$ZV(Y)(f(3)+o(1))$ and $ZV(Y)(F(3)+o(1))$.
The error is uniform as $Z$ stays between fixed powers of $X$;
\eqref{eq:BV} pays all progression remainders.  Substituting
\eqref{eq:prod} and $Z=X/m$ yields \eqref{eq:permbound}.
\end{proof}

\begin{proof}[Proof of Theorems \ref{th:shifted} and \ref{th:rough}]
First, the exact gate satisfies
\[
 0\le G_2(mp+2)\one_{\{mp+2\in\aset_X\}}
 \le\one_{\{\phimin(mp+2)>Y\}}
\]
by Lemma \ref{lem:G}.  The upper half of Lemma \ref{lem:sieve}
therefore applies to $S_{+,3}^{\Box}$ after summing over the
allowed \emph{true} primes $r,s$.  Its lower and upper halves apply
to $\mathcal R_X^{\Box}$.

By the prime reciprocal-sum form of Mertens' theorem,
\[
 \sum_{X^{a_1}<r\le X^{a_2}}\frac1r
  =\log\frac{a_2}{a_1}+o(1),\qquad
 \sum_{X^{b_1}<s\le X^{b_2}}\frac1s
  =\log\frac{b_2}{b_1}+o(1).
\]
The exponent assumptions give $r<s$ and $rs>D$ automatically,
with no representation multiplicity.  Since
$B_X\sim2C_2X/\log X$ and $\log Y=(\log X)/21$,
dividing the summed sieve upper by $B_X$ yields
\[
 \frac{(4C_2/3)X}{B_X\log Y}
       \longrightarrow14.
\]
The corresponding lower has the extra $\log2$ factor.
This proves \eqref{eq:boxbound} and \eqref{eq:rough}.
For the entire positive sector, the same prime harmonic limit as
in \eqref{eq:primeharmonic} gives $\mathcal I$, so its upper is
$14\mathcal I$.  Substitution of the four rational endpoints in
\eqref{eq:box} gives \eqref{eq:numericbox}.
The original rank and squarefree gates on $mp+2$ are retained
until the displayed one-sided inequality; they have not been
replaced by a conjectural prime-pair main term.
\end{proof}

\section{Why an asymptotic sieve does not follow from this input}

A stronger theorem on primes in thin sequences is the
Friedlander--Iwaniec asymptotic sieve \cite{FI1998}.  Its
bilinear hypothesis (B), besides its independent remainder
hypothesis, is a genuine parity-breaking condition.
Here is an exact obstruction to applying that bilinear hypothesis
\emph{without new information} to the three-prime input.

\begin{proposition}[A failed bilinear input on a nonempty rough family]
\label{prop:fi}
For all sufficiently large $X$, define a nonnegative sequence $a_M$
supported on $M\in(X,2X]$ having the representation
\[
 M=rsp,\quad X^{0.127}<r\le1.05X^{0.127},\quad
 X^{0.201}<s\le1.05X^{0.201},\quad r,s,p\text{ prime},
\]
and set $a_M=(\log p)\one_{\{\phimin(M+2)>Y\}}$ there, zero elsewhere.
Then $\sum_Ma_M\gg X/(\log X)^3$.
For the Friedlander--Iwaniec bilinear expression with parameter
$C=1$ and the dyadic divisor range $H<v\le2H$ where $H=X^{0.328}$,
one has the \emph{exact} equality
\begin{equation}\label{eq:FI}
 \sum_u\left|\sum_{\substack{H<v\le2H\\uv\le2X}}
  \gamma(v,1)\mu(uv)a_{uv}\right|=\sum_Ma_M,
 \qquad \gamma(v,C)=\sum_{d\mid v,d\le C}\mu(d).
\end{equation}
Thus even a single fixed positive power of logarithmic saving over
the total mass is impossible for this sequence indexed by \(M\).
Applying the asymptotic sieve to \(N=M+2\) would involve the distinct
sequence indexed by \(N\); equation~\eqref{eq:FI} concerns the \(M\) indexing.
This assertion concerns the larger rough family; no positive mass
is thereby established for the exact gate $M+2\in\aset_X$.
\end{proposition}

\begin{proof}
The lower linear sieve in Lemma \ref{lem:sieve} applies uniformly
to the given $r,s$.  By the prime number theorem in fixed-ratio
intervals,
\[
 \sum_{X^a<q\le1.05X^a}\frac1q
  \sim\frac{\log1.05}{a\log X}
 \qquad(a=0.127,0.201).
\]
The product of these two sums and the sieve lower is a positive
constant times $X/(\log X)^3$.
Every integer on the support has exactly the three distinct prime
factors $r,s,p$, so $\mu(M)=-1$.  Its divisors in
$(X^{0.328},2X^{0.328}]$ consist of the unique divisor $v=rs$:
$r,s$ are smaller by fixed powers of $X$, and $p,pr,ps,M$ are
larger by fixed powers.  Further, $\gamma(v,1)=1$.
Thus each inner sum in \eqref{eq:FI} is a sum of negative
numbers, and upon removing the absolute values and changing sign,
every $a_M$ is counted exactly once.  This proves the equality.
For instance, taking the asymptotic-sieve parameters
$D_{\mathrm{FI}}=(2X)^{0.70}$ and
$\Delta=(2X)^{0.04}$ places $H\asymp X^{0.328}$ strictly between
$\sqrt{D_{\mathrm{FI}}}/\Delta$ and
$\sqrt{2X}/(\log X)^{10}$, so this is a genuine tested bilinear
range, not a divisor range outside the theorem's hypotheses.
No claim is made that the remaining, unrelated hypotheses of
\cite{FI1998} hold for the sequence $a_M$.
\end{proof}

\section{Normalization, comparison, and outstanding problem}

Theorem \ref{th:unshifted} is a \emph{one-site} result: it
counts $rsp$ without testing $rsp+2$.  In contrast,
Theorem \ref{th:shifted} is about the \emph{actual} pair
$(rsp,rsp+2)$ and keeps the squarefree, factor-count, and large-factor
conditions on its right member, but furnishes only an upper bound.
Theorem \ref{th:rough} proves positive shifted mass only after replacing
those right-member requirements by roughness.
The conditions $r,s,p$ prime do not make $rsp+2$ prime or a
product of two or three primes.

For clarity, the containing signed quantity studied in the originating
research is
\begin{equation}\label{eq:fullsigned}
 R_X=\frac1{B_X}\sum_{\substack{M\in(X,2X]\ \mathrm{squarefree}\\
                   2\le\omega(M)\le7,\ \phimin(M)>X^\delta\\
                   M+2\in\aset_X}}
               G_2(M+2)F_D(M).
\end{equation}
The rank-two part is nonpositive by Theorem \ref{th:classification}
and Lemma \ref{lem:G}.
The positive rank-three part is \(S_{+,3}(X)/B_X\); higher ranks may have
additional positive and negative terms.  Hence
\eqref{eq:wholebound} \emph{cannot} be used as an upper bound for
$R_X$ by dropping all higher ranks without first establishing their
signs or separately bounding their positive parts.  Nor can the
negative rank-two contribution be counted as an independent saving
without matching it to the same complete summation.

The local estimate \eqref{eq:numericbox} leaves the complementary
boxes uncontrolled. A positive lower bound for genuine twin primes
remains open.

For context, recent dispersion theorems impose specific
factorizations or modulus weights, rather than accepting the
arbitrary joint rank gate $G_2(rsp+2)$ as a product of independent
coefficients; see \cite{Maynard, PascadiSmooth, PascadiLevels}.
An application to \eqref{eq:fullsigned} would require a new
arithmetic estimate keeping the fixed shift, all relevant factor
ranks, and the additive precision of the signed sum.

\paragraph{Provenance and verification.}
This standalone note is extracted from an internally developed
twin-prime research record, notably its sections 308LQ--308LU and
the additional sign and shifted-sieve calculations dated
19--20 November 2026.  The elementary formulas and the constants
in this note have been reconstructed for this manuscript.
The analytic sieve is invoked only in its standard established range.
The note has not undergone external peer review or a formal proof check.
Substantive initial research and manuscript preparation involved
OpenAI ChatGPT under human direction.  Authorship and the correspondence
address are editable fields at the top of this \LaTeX{} source.

\appendix
\section{Numerical consistency check for the dilogarithm constant}
The following optional Python 3 computation uses \texttt{mpmath}.
It is not a substitute for the integration proof in
Theorem \ref{th:unshifted}.
\begin{verbatim}
import mpmath as mp
mp.mp.dps = 40
delta, t = mp.mpf('0.05'), mp.mpf('0.24995')
a = delta / t
I = (mp.polylog(2, 1-a) - mp.polylog(2, a)
     + mp.log(1/a)*mp.log(1/(1-a))) / 2
J = mp.quad(lambda u: mp.log(t/max(delta, t-u))/u,
            [delta, t-delta, t]) / 2
box = (14 * mp.log(mp.mpf('0.1305')/mp.mpf('0.125'))
       * mp.log(mp.mpf('0.202')/mp.mpf('0.200')))
print(I, J, I-J, box, 14*I)
\end{verbatim}

\begin{thebibliography}{99}
\bibitem{FI1998}
 J. Friedlander and H. Iwaniec,
 \emph{Asymptotic sieve for primes},
 Ann. of Math. (2) \textbf{148} (1998), 1041--1065.
 \url{https://annals.math.princeton.edu/articles/13036}.
\bibitem{OperaDeCribro}
 J. Friedlander and H. Iwaniec,
 \emph{Opera de Cribro}, Amer. Math. Soc. Colloquium Publ.,
 vol.~57, 2010.
\bibitem{MZA}
 K. Matom\"aki and S. Z\'u\~niga-Alterman,
 \emph{Weighted sieves with switching},
 Math. Proc. Cambridge Philos. Soc. \textbf{179} (2025), 351--372,
 Lemma~2.5 and the linear-sieve functions in \S2.
 \url{https://doi.org/10.1017/S0305004125000258}.
\bibitem{Maynard}
 J. Maynard,
 \emph{Primes in arithmetic progressions to large moduli I:
 Fixed residue classes},
 \url{https://arxiv.org/abs/2006.06572}.
\bibitem{PascadiSmooth}
 A. Pascadi,
 \emph{Smooth numbers in arithmetic progressions to large moduli},
 Compos. Math. \textbf{161} (2025), 1923--1974.
 \url{https://doi.org/10.1112/S0010437X2500747X}.
\bibitem{PascadiLevels}
 A. Pascadi,
 \emph{On the exponents of distribution of primes and smooth numbers},
 \url{https://arxiv.org/abs/2505.00653}.
\bibitem{BV}
 E. Bombieri,
 \emph{On the large sieve}, Mathematika \textbf{12} (1965), 201--225;
 A. I. Vinogradov, \emph{The density hypothesis for Dirichlet
 $L$-series}, Izv. Akad. Nauk SSSR Ser. Mat. \textbf{29} (1965),
 903--934.  (Classical Bombieri--Vinogradov theorem.)
\end{thebibliography}
\end{document}
